\subsection{Transport resolvents and preservation of memory}
\label{subsec:ii-memoria-resolvente}

The prolonged history determines the windows
\(I_m=\{9m+1,\ldots,9m+9\}\), for \(m\geq0\).
With the numbering of this section, \(I_m=E_m\) during the first cycle.
Let \(\sigma_m\in\{-1,1\}\) be the orientation preserved by that history;
for \(0\leq m<12\), it coincides with \(\Sigma_{12}(m)\). Define
\[
 a_m=\sigma_m\sum_{t\in I_m}\mathbf1_{\{2,4,6,8\}}(d_t),
 \qquad |a_m|\leq9.
\]
The record and the moving frame are prolonged by
\[
 \mathbf z_{m+1}=\mathbf z_m+a_m\mathbf e_{m\bmod12},\qquad y_m=S^{-m}\mathbf z_m.
\]
The initial state \(\mathbf z_0\) preserves the preceding history; it is zero for
an initially empty prefix. In this subsection we write
\(R=S^{-1}\) for frame transport. The update
\eqref{eq:eta} is retained in every window. Subsequent orientation and events
are obtained from the prolonged history: frame return
does not presuppose periodicity of the event sequence.

\begin{proposition}[Truncation identity with terminal state]
For $N\geq1$, let
\[
Y_N(u)=\sum_{m=0}^N y_mu^m,\qquad
H_{\eta,N}(u)=\sum_{m=0}^{N-1}\eta_mu^m.
\]
The following exact polynomial identity holds:
\begin{equation}
(I-uR)Y_N(u)=y_0+uH_{\eta,N}(u)-u^{N+1}Ry_N.
\label{mem:identidad-finita}
\end{equation}
\end{proposition}
\begin{proof}
In degrees $1\leq m\leq N$, the coefficient on the left-hand side is
$y_m-Ry_{m-1}=\eta_{m-1}$, by \eqref{eq:eta}.
Its coefficients in degrees zero and $N+1$ are $y_0$ and $-Ry_N$,
respectively. Comparing coefficients proves the equality.
\end{proof}

The terminal term preserves the transported memory when a history is cut.
The identity allows a finite evaluation to be compared with its prolongation
while keeping explicit the state left at the boundary.

\subsubsection{Correspondence with common-core transport}

The window numbering of Section~\ref{subsec:memoria-resolvente}
and that used here describe the same partition of the history. The
correspondence is
\[
 \underbrace{I_m=E_{m+1}}_{\text{common-core index}}
 \quad\longleftrightarrow\quad
 \underbrace{I_m=E_m}_{\text{index in this section}},
 \qquad 0\leq m<12.
\]
The subscript changes the name of the window; its nine transitions,
its orientation \(\sigma_m=\Sigma_{12}(m)\), its event \(a_m\), and
its position in the record remain the same. For every prolongation,
\(R=S^{-1}\), \(y_m=S^{-m}\mathbf z_m\), and
update~\eqref{eq:eta} coincide with
\eqref{mem:actualizacion}. Orientation and the event sequence
retain their dependence on the history, also after
frame return.

This identification allows Proposition~\ref{mem:prop-serie}
to be applied directly: series~\eqref{mem:serie},
bound~\eqref{mem:cota}, and uniform convergence of its derivatives
correspond to this same record and the same window clock. The
finite identity~\eqref{mem:identidad-finita} also determines the terminal
term when that series is truncated. Its hypotheses retain the initial history
and the bounded event \(|a_m|\leq9\).

When the record is reused, Proposition~\ref{mem:prop-schur} jointly
preserves the operator, source, and reader through
\eqref{mem:schur-fuente} and~\eqref{mem:schur-lector}. Example
\eqref{mem:ejemplo-ciclo} specifies the effect of the excursions omitted
by a compression. The affine composition~\eqref{mem:cociclo-tres}, the
intermediate blocks~\eqref{mem:schur-intermedio}, and identity
\eqref{mem:schur-transitivo} ensure consistency when segments
are concatenated or several layers eliminated. Their complete proofs reside
in Section~\ref{subsec:memoria-resolvente}, with the domains,
sources, and reading maps specified there.

Balance~\eqref{ii:mem:balance} preserves the correspondence between the
oriented readers of the signed record and the common quadratic form.
The distinction between the accumulated vector \(\mathbf z_m\) and the
incidence decimal block \(z_m\) remains that established in this section.
The transport identities apply to memory thus typed; they do not
by themselves determine additional analytic coefficients.
